%
%  Conclusion
%
\chapter{Conclusion}
\label{ch_6}
This thesis provides an overview of wave-mechanics as applied to LSS. The interpretation of the Schr\"odinger equation was re-examined in Chapter \ref{ch_3}, previous inconsistencies and misconceptions were explored and resolved to the best of our ability. Our understanding of wave-mechanics as applied to LSS should now be clearer and more concise. The approach in this thesis is purely classical due to the high particle occupation number in each quantum state, so pure quantum mechanics is never present. The quantum nature of dark matter particles is mostly unknown but not thought to be significant. However this does not affect the main outcome of this work.

Chapter \ref{ch_3} made clear the reason for choosing the Schr\"odinger equation, and not a generic wave-equation, for simulating LSS using a wave-mechanical method. The Schr\"odinger equation is an energy equation where the terms can be easily interpreted in a physical manner using well-developed techniques; there is an obvious link between observables and operators; and the wavefunction is a single complex field that provides a simple method of obtaining the density and velocity.

We have shown the correct interpretation of the Schr\"odinger in this context is that of a generic energy equation that describes the evolution of waves but does not necessarily have to describe a quantum system. In a classical system the effective Planck constant is now analogous to the diffusion coefficient of fluid dynamics. In order to describe such a classical system of $N$ particles one must `multiply up' the single-particle Schr\"odinger equation which assumes a high occupation number.

Also included in Chapter \ref{ch_3} is a modern derivation of the fluid equations from the Schr\"odinger equation. This is lacking in previous  literature. In addition to providing a sketch of the derivation we delved further into the interpretation of the equations, building upon a translation of the original publication by Madelung (see Appendix \ref{app_m}), as well as new insights from later literature and further study.

Crucially, we believe that it is important not to confuse the role of the so-called pressure term. It is a term that highlights the difference between a free particle and a fluid. It is not a term that is `added' in by the Schr\"odinger equation; there is no new or added pressure into the system. In fact, the opposite is closer to the truth. The free particle Schr\"odinger equation describes a fluid that has no internal pressure. Short discovered (under the limit of $\hbar \rightarrow 0$) that the term is only important in regions of high-density. This suggests that the subtraction of the pressure term is negligible and hence a free-particle is the same as a fluid in regions of low density. In high density regions it is natural to expect the pressure to dominate the fluid's evolution.

Johnston, considered a different approach from myself and Short and adds in (by hand) the pressure term in her calculations. Thus, the wavefunction that she describes is truly a fluid.

In Chapter \ref{ch_4}, we presented the results of our investigation of the Free Particle Approximation (FPA) model and confirmed the robustness of the results provided by Short. Our tests were based upon the mathematics provided in his thesis but were carried out using entirely independent codes -- thus representing a fully independent test that reinforces the idea that the FPA is a sound approximation scheme. We were able to conclude that the FPA is a fast and efficient method for probing the quasi-linear regime of density perturbations, it proves to be a good match for the Zel'dovich approximation and breaks down at shell crossing as the Zel'dovich method does. Short proved that the FPA is formally equivalent the Zel'dovich approximation with adhesion (in the limit of $\nu \rightarrow 0$).

The 1D and 3D toy models considered in Chapter \ref{ch_4} were a demonstration of a consistent framework. I have shown that the FPA can handle cosmological initial conditions by providing results that are comparable to the widely available $N$-body codes (at least in the quasi-linear regime). The benefit of the FPA is that it runs much faster than all other known methods of simulation. It merely requires a single step of computation due to the symmetry of the equations.

After developing and testing the FPA, I went on to investigate another method of solving the Schr\"odinger equation, colloquially we call it the Cayley method (Chapter \ref{ch_5}) as it involves describing the exponentials via Cayley's approximation. In this approach we no longer assume the trick that was inherent to the FPA, where the effective potential is zero. In the Cayley approach we solved the full system of equations: both the kinetic term and solved Poisson's equation of gravity.

As shown in Chapter \ref{ch_5}, this process is trickier than it may first appear. Under the recommendation of Watanabe's publication we adopted Suzuki splitting operators and also adopted Widrow \& Kaiser's method of re-writing the equations into a simplified form that deals with timesteps as equal steps in $\ln a$. Fortunately, each component works consistently with one another. The simulations prove to be robust and stable. Consequently, the main goals of developing a wave-mechanics code for LSS were achieved. Those goals were: (1) 3D coordinates, (2) self-consistent gravity, (3) expanding coordinates, (4) periodic boundaries, (5) mass conserving (as a bonus it conserves momentum too). In addition this also satisfies the original goals suggested by Widrow \& Kaiser: A model that describes collisionless matter; matter is described as a field, not particles: that is, continuous; function only of space and time (3+1 d); follow multiple streams in phase space (`hot' / dispersive); competitive with $N$-body techniques for computing time. Sabiu's run with Gadget took approximately the same time to complete as my Wave-mechanics code, of the order of a few hours; however, no rigorous speed testing was conducted in this thesis. Both codes were run on different machines so does not represent a fair comparison but I do expect both codes to be a comparable speed. The number of operations required in each case are similar in magnitude.

The last of our goals (mass conserving) is an indicator of the reliability of our results. Not only did we achieve our goals but have done so in a manner that is consistent with the physics. This was only possible by our strict choice of a symplectic integrator such that the norm of the wavefunction is conserved over time. Hence, this choice of integrator is consistent with unitary time evolution in Quantum Mechanics and further hence that it should be less of a surprise that momentum is conserved as well as mass.

Through out this work I have highlighted the similarities between wave-mechanics and $N$-body simulations. The latter was used as a benchmark for speed and accuracy. However, I have shown wave-mechanics provides a different method for simulating Large Scale Structure. It has many benefits but it lacks the maturity and development seen in modern $N$-body codes. It seems fair to assume that if wave-mechanics is developed further and allowed to run on high-performance machines then it should have no difficulty matching and, perhaps, beating equivalent $N$-body codes for speed and accuracy (in terms of density resolution). This is a guess based on the assumption that the wave-mechanics code can be written to have fewer or less expensive operations than an $N$-body code. In terms of resolution, we believe that a continuous density field representation should ultimately lead to better resolution of the density field. We also believe that wave-mechanics should provide a better representation of the Universe (over $N$-body codes) because they describe a continuous density field and suffers less from discretization problems. The allowance of multi-streaming also prevents two-body relaxation that is a critical issue for $N$-body codes, the current solution to that problem is unphysical and is completely avoided in wave-mechanics.

The main problems with the results presented in this thesis were the unexpected, although not necessarily inconsistent, results for density and velocity. The messiness of the evolved density fields is potentially due to interference, this is unique to a wave-mechanics but not something that desire from a classical system. This effect requires further investigation.

The simulations presented were run at a coarse spatial resolution of $64^3$, this was a problem for generating our initial conditions and hence our subsequent results. At higher resolutions the results look better behaved but were prohibitively more expensive to run within the time frame of this thesis. In order to make sense of the lower resolution simulations we smoothed the data using a gaussian filter (essentially a low pass filter that takes out high frequency `noise'). This operation gave data that looked closer to the Gadget data. We would have liked to use the raw, unsmoothed data, for our point of comparison but it did not look consistent with what we first expected.

The histograms of both density and velocity are initially gaussian and this distribution is fairly well maintained for velocity at all times, while the density becomes a skewed gaussian. The under-dense regions in our wave-mechanics results are evacuated quicker than seen in an $N$-body code but the peaks do not seem to collapse as fast, nor do the peaks reach the same height as seen in the results from Gadget. The choice of the $\nu$ parameter influences the dynamics in such a way that it can greatly affect the collapse speed; in some cases, collapse can be completely prevented. Naturally, the lack of evolution would be inconsistent with what we observe in the real Universe.

These issues will require further investigation but it is our belief that they do not rule out wave-mechanics as a viable method for simulating Large Scale Structure. One of the main concerns we have had is the lack of time needed to run high resolution simulations, although we did some very simple tests at $128^3$ our final runs were at $64^3$ resolution. Initially, we incurred problems from running low resolution simulations where collapse was too fast in the initial timesteps. As mentioned in section \ref{cos_sim}, this effect disappeared for simulations at a resolution of $128^3$. The effect also went away whenever we applied gaussian smoothing to our initial density field.

Given the positive results from this work, we now propose future directions and possible extensions to wave-mechanics codes. We would greatly welcome any reader to reproduce the results from this thesis and / or attempt to implement any of the future work.

\section{Future work: extending the full Schr\"odinger-Poisson system}
In this final section I present ideas that expand upon the standard techniques of wave-mechanics as presented in the previous chapters of this thesis. The aim of this section is for me to develop ideas that will help advance wave-mechanics further and show that our method of evolving structure formation is capable of achieving everything that can be done in an $N$-body code. I see my thesis and all previous work as being the first generation of wave-mechanics codes. The next generation of codes will build upon all previous literature so that they are more diverse in what they achieve, while being also robust.

The best $N$-body codes (such as GADGET and RAMSES) have many more features than current wave-mechanical codes presented hitherto. A short list of features that I believe are desirable for the next generation of wave-mechanics codes, and are already prominent in the best $N$-body codes, are: 
(1) multiple particle/fluid species, (2) adaptive mesh refinement, (3) parallelization.

In this chapter I will review each item in the above list and expand upon how one might implement these features into a wave-mechanics code. Fortunately these ideas have already been explored to some extent in previous publications; however, they are yet to be featured in a full cosmological wave-mechanics code. Johnston \citep{bib_rj} has already written a code that models two species of fluid and Watanabe \citep{bib_wat} has provided many ideas that can be included in a next generation of wave-mechanical LSS codes: such as a possible method for implementing mesh refinement and parallelization. The following ideas are presented in the remainder of this thesis:

\begin{itemize}
\item Multiple fluids
\item The splitting operator approach to including more physics (such as pressure)
\item Periodic boundaries via adhesive operators
\item Mesh refinement
\item Parallelization
\item Including vorticity and spin
\end{itemize}

\noindent The last bullet point will be presented in the next chapter, they are less conventional and hence far more speculative. They are effects that are not currently accounted for in standard LSS simulations. These ideas are the inclusion of gravitomagnetism and spinning objects into a wave-mechanics code. Both ideas might allow wave-mechanics codes to potentially probe beyond the standard model of cosmology.


\subsection{Multiple fluids}
The idea of a multiple particle-species is not new to $N$-body codes. Special variants \citep{bib_hmpc1} of the Hydra code \citep{bib_hmpc} allows for Baryons as well as CDM particles to be present, these codes have been written to account for gas dynamics where the Baryons will behave as dispersive particles with an associated temperature.

Johnston has shown how to include another fluid of the same mass into a Schr\"odinger code. The two species have separate wavefunctions with separate Schr\"odinger equations except that they have a common gravitational potential. In principle this is easy to extend to $N$ species of particles (or fluids). The initial mass of each fluid is allowed to be different, this will be expressed by the fact that the total mass of the wavefunction for each species will be different. It would also be possible to alter the value of $\nu = \hbar/m$ for each species, hence account for different dispersion properties.

I believe that the splitting operators presented in this thesis provide a natural way of adding new particle species in a modular way. Here I will briefly recapitulate how this method evolves the wavefunction:

\begin{equation}
\psi(x, t+dt) = e^{-i (K+V) dt} \psi(x, t) = e^{-i K dt/2}e^{-i V dt}e^{-i K dt/2} \psi(x, t)
\end{equation}

\noindent We have chosen to split the kinetic energy operator such that one half timestep is performed before and then after the potential energy operator. I envisage performing each half-step kinetic operator sequentially for 1 to $N$ (the number of species / fluids) then performing the calculation of a common gravitational potential. Here is an overview of how I expect the evolution part of the algorithm to look:

\begin{enumerate}
\item $\psi_{1 \ldots N}(t + \frac{1}{2} dt) = \hat{K}_{1 \ldots N} \; \psi_{1 \ldots N} (t)$
\item $\psi_{1 \ldots N}(t + \frac{1}{2} dt) = \hat{V}_{1 \ldots N} \; \psi_{1 \ldots N} (t + \frac{1}{2} dt)$
\item $\psi_{1 \ldots N}(t + dt) = \hat{K}_{1 \ldots N} \; \psi_{1 \ldots N} (t + \frac{1}{2} dt)$
\end{enumerate}

\noindent Here I denote the kinetic energy operator by $\hat{K}$ and the potential energy operator by $\hat{V}$. The subscripts ${1 \ldots N}$ denote the different particle species.

The operator $\hat{K}$ above is stated generically so the kinetic operator for each species could have a different value of $\nu$. It would also be possible to modify the kinetic energy operator of each species too, although there are limits to how far the operator can be modified: for example, the addition of any term that does not commute with the kinetic energy is very likely to break the energy conservation of the code.

To properly account for additional physics one would need to split the operators (in a nested fashion) as Suzuki suggests. For example it should be possible to include Baryons having internal interactions.

\subsection{Including additional physics}
\label{addphys}
If I was to add pressure into my code (as Johnston did) I think I would need to include it as a separate split-operator. I don't believe that it would be entirely correct to modify the kinetic energy operator to account for pressure, for example simply writing $\hat{\tilde{K}} = \hat{K} + \hat{P}$. The evolution would then be written as (this is in short hand form, I omit $i, \; {\rm dt},$ etc in the exponential):

\begin{equation}
\exp(\hat{\tilde{K}})\psi(t) = \exp(\hat{K} + \hat{P}) \psi(t)
\end{equation}

\noindent To properly account for pressure in a robust manner in the Cayley method then one would need to separate the operators in a way that is consistent with Suzuki's method of decomposition. The full decomposition of the Hamiltonian (including all 3 terms: kinetic, potential and pressure) is:

\begin{equation}
\exp(\hat{K}+\hat{V}+\hat{P})\psi(t)\rightarrow  \exp(\frac{1}{2}\hat{K})\exp(\frac{1}{2}\hat{P})\exp(\hat{V}) \exp(\frac{1}{2}\hat{P})\exp(\frac{1}{2}\hat{K}) \psi(t)
\end{equation}

\noindent I believe this would be true for the inclusion of any number of operators, provided they are split up in the manner that Suzuki suggests. In a second paper from Watanabe \& Tsukada \citep{bib_wat2}, they show how to use the Suzuki splitting operators to include a magnetic potential when simulating the dynamics of an electron in a magnetic field.

\subsection{Periodic boundaries via adhesive operators}
As mentioned in section \ref{per_bound}, Watanabe \citep{bib_wat} suggested a method for implementing periodic boundary conditions that he coined `adhesive operators' (not to be confused with the Zel'dovich adhesion model). These `operators' connected one boundary to another (connects the boundaries like an adhesive) such that waves that exited on the left-hand-side of the system and would reappear on the right-hand-side. I believe that if these operators can be successfully implemented in a cosmological wave-mechanics code then it would save on computing time as my method (section \ref{per_bound}) involves a double iteration of each recursion relation can be avoided. As already mentioned, I made a brief attempt to implement periodic boundaries in this way but it did not conserve mass. Any possible method that can cut computing time but preserve unitarity would be worth investigating.

Watanabe's idea is simple and elegant; the mathematics is relatively straight forward but it seems that not enough details are provided in the publication of Watanabe to ensure perfect implementation. Also, as previously mentioned, it is not obvious that such operators would overcome the inherent problem of the recursion relations for the auxiliary functions ($e$ and $f$).

Watanabe's adhesive operators are ubiquitous, he uses them for the implementation of periodic boundaries, mesh refinement and parallelization. Hence, we believe that such a method if it can be independently proven to work would yield a powerful method for improving speed, accuracy and size of simulations.

Here is a brief sketch of Watanabe's idea. He notes that in a crystal the wavefunction can be taken as periodic, that is to say that some region far from the real boundaries of the crystal will look like a repeating unit of a larger lattice. Hence, the extent of the lattice is pseudo-infinite. This is the same type of argument that is made for using periodic boundary conditions in cosmological simulations. The wavefunction obeys the periodic relation:

\begin{equation}
\psi({\bf r} + {\bf R},t )= \psi({\bf r},t)e^{(i\phi)}, \;\;\; \phi = {\bf k}.{\bf R}
\end{equation}

\noindent ${\bf k}$ is the Bloch wavenumber and ${\bf R}$ is the length of the lattice (although Watanabe calls this the unit vector, implying it should be the spacing between gridpoints). Bloch's theorem requires the modulus of the exponential to be one, hence the modulus of the wavefunction is also preserved. This is necessary for a smooth transition at the boundary and for unitary time evolution. For code implementation, these modifications appear in the $\Omega$ function of Goldberg. Recall that when using splitting operators the potential $V$ is separate from the kinetic operator $K$ and hence taken out of the $\Omega$ function as seen in Goldberg. Hence the adhesive operators are a modification to the kinetic energy operator. Copying Watanabe's notation we write the matrices for the time evolution in the following manner ($V$ is omitted):

%\[
\begin{eqnarray}
\left(\begin{array}{cccccc}
A & -1 & 0  & 0 &\ldots & e^{+i\phi} \\
-1 & A & -1 & 0 &\ldots & 0 \\
0 & -1 & A  & -1&\ldots & 0 \\
\vdots & & & & & \vdots \\
e^{-i\phi}& 0 & \ldots & 0 & -1 & A
\end{array}\right)
\left(\begin{array}{c}
\psi(0,t + dt) \\
\psi(1,t + dt) \\
\psi(2,t + dt) \\
\vdots \\
\psi(L-1,t + dt) \\
\end{array}\right) \\
=
\left(\begin{array}{cccccc}
B & -1 & 0  & 0 &\ldots & e^{+i\phi} \\
-1 & B & -1 & 0 &\ldots & 0 \\
0 & -1 & B  & -1&\ldots & 0 \\
\vdots & & & & & \vdots \\
e^{-i\phi}& 0 & \ldots & 0 & -1 & B
\end{array}\right)
\left(\begin{array}{c}
\psi(0,t) \\
\psi(1,t) \\
\psi(2,t) \\
\vdots \\
\psi(L-1,t) \\
\end{array}\right)
\end{eqnarray}
%\]

\noindent the definitions of $A$ and $B$ are not necessary for understanding the method. In the simplest case $A$ and $B$ take the expected form as you would expect from the kinetic energy matrices as they appear in section \ref{onedim}. Here we use $L$ to denote the length of the lattice, so $L-1$ is the last element of the array.

As Watanabe notes, the addition of these off-diagonal elements should prevent the matrices finding an efficient solution. The usual methods of inversion work well for diagonal, or band-diagonal, matrices but not non-diagonal matrices. However, the trick is to note that the first line and the last line of the operator matrices can be solved independently from the rest of the matrix: $\partial^2_x = \partial^2_{x-td} + \partial^2_{x-ad}$. The kinetic energy can be written as the matrix addition of the tridiagonal contribution $\partial^2_{x-td}$ and the off-diagonal terms (the adhesive operator) $\partial^2_{x-ad}$.

An alternative method for periodic boundaries would be the use of Fourier transforms as they are inherently periodic. However, our attempts to implement the kinetic energy operator via an FFT method have been unsuccessful.

\subsection{Mesh refinement}
Mesh refinement is a method of improving the code's efficiency. Regions that are known to grow into large densities can be pre-binned before the simulation starts. That is, the mesh can be of a finer resolution in areas of higher density. Alternatively, a more computationally expensive way would be to allow the computer to determine how to bin the data on the fly using so-called Adaptive Mesh Refinement (AMR). Static mesh refinement requires the user to know where the high density regions form and would require a lot of effort for each new simulation.

AMR is not a new idea, nor are ideas of adaptive quadrature, but such an idea is missing in the current wave-mechanics of LSS literature.
Watanabe (\citeyear{bib_wat}) has shown how to include mesh refinement in a Schr\"odinger solver. Mesh refinement is easily achievable for the kinetic energy operator by use of Watanabe's adhesive operators. The potential energy calculation requires separate consideration. Such an idea is particularly attractive for LSS simulations and should perhaps be one of the first avenues for future wave-mechanics research. It can provide increased spatial resolution in the areas that need it most.

For mesh refinement, Watanabe decomposes the kinetic energy operator into two parts: one part that is easy to solve (the diagonal) and the part that requires a separate calculation in the manner of the adhesive operator. The main difference between the refined calculation and the non-refined is that the kinetic energy operator may require more than 3 gridpoints for calculation and that the weightings for each gridpoint are not necessarily the usual -1, +2, -1 as seen for the central difference method. The weightings can be fractional, as in equations (32 - 34) of Watanabe (\citeyear{bib_wat}).

\subsection{Parallelization}
Another key development for increasing the size of simulations is to develop a way of parallelizing the wave-mechanics code. Some parallel Schr\"odinger solvers \citep{bib_yee, bib_schn, bib_str} exist but their method is different from the Goldberg/Cayley scheme as used in this thesis. Furthermore, it is not entirely obvious how reliable these codes are. The recursion relations are inherently difficult, if not impossible, to put into a parallel code. Ideally, the wavefunction can be split up into independent regions and then evolved by $dt$. The regions would then be reconciled with the whole simulation at the end of each time step.

For parallelization Watanabe suggests the use of his adhesive operators again: they can connect two regions, or two boundaries, with ease. Each region is calculated independently and then matched at the boundary.

Without the adhesive operators, one would need to run a single set of recursion relations over multiple processors. Each processor would rely upon the processor before it in the chain. This is much slower than if each processor can work independently then pass data at the end of each calculation set. This becomes very expensive when periodic boundaries are required.

In the absence of periodic boundaries, where only one recursion is necessary, the recursion relations can be staggered: while the second processor continues the recursion relation for $e$, processor one could start the recursion relation for $f$. When periodic boundaries are implemented this inevitably means that the first processor is idle while the next second processor continues with the recursion. This is necessary as the second iteration of each recursion relation requires the last value in the sequence (comes from the last processor) before starting again.

The lack of a clear parallelization scheme will inhibit progress in wave-mechanics. Future research into this area may have to choose a different Schr\"odinger solver if such a solver is able to be parallelized while being reliable.



